% 第 2 周周记（2026.08.10 - 08.14）
\documentclass[UTF8,fontset=fandol]{ctexart}
% =====================================================================
% internship/preamble.tex
% 实习周记统一模板的共用导言区。各周文件在 \documentclass 之后用
%   % =====================================================================
% internship/preamble.tex
% 实习周记统一模板的共用导言区。各周文件在 \documentclass 之后用
%   % =====================================================================
% internship/preamble.tex
% 实习周记统一模板的共用导言区。各周文件在 \documentclass 之后用
%   \input{../preamble.tex}
% 引入本文件，正文前用 \weekhead{周次}{时间}{主题} 输出抬头。
% 编译命令：xelatex weekN.tex
% =====================================================================

\usepackage[a4paper,top=2.6cm,bottom=2.6cm,left=2.8cm,right=2.8cm]{geometry}
\usepackage{booktabs}   % 表格横线，week3 的结果对比表会用到
\usepackage{array}

% 周记抬头：\weekhead{周次}{时间区间}{本周主题}
\newcommand{\weekhead}[3]{%
\begin{center}
{\zihao{3}\bfseries 实习周记（#1）}\\[10pt]
{\zihao{-4}姓名：【姓名】\qquad 学号：【学号】}\\[4pt]
{\zihao{-4}时间：#2}\\[4pt]
{\zihao{-4}主题：#3}
\end{center}
\vspace{2pt}
\hrule height 0.8pt
\vspace{10pt}
\zihao{-4}%
}

% 引入本文件，正文前用 \weekhead{周次}{时间}{主题} 输出抬头。
% 编译命令：xelatex weekN.tex
% =====================================================================

\usepackage[a4paper,top=2.6cm,bottom=2.6cm,left=2.8cm,right=2.8cm]{geometry}
\usepackage{booktabs}   % 表格横线，week3 的结果对比表会用到
\usepackage{array}

% 周记抬头：\weekhead{周次}{时间区间}{本周主题}
\newcommand{\weekhead}[3]{%
\begin{center}
{\zihao{3}\bfseries 实习周记（#1）}\\[10pt]
{\zihao{-4}姓名：【姓名】\qquad 学号：【学号】}\\[4pt]
{\zihao{-4}时间：#2}\\[4pt]
{\zihao{-4}主题：#3}
\end{center}
\vspace{2pt}
\hrule height 0.8pt
\vspace{10pt}
\zihao{-4}%
}

% 引入本文件，正文前用 \weekhead{周次}{时间}{主题} 输出抬头。
% 编译命令：xelatex weekN.tex
% =====================================================================

\usepackage[a4paper,top=2.6cm,bottom=2.6cm,left=2.8cm,right=2.8cm]{geometry}
\usepackage{booktabs}   % 表格横线，week3 的结果对比表会用到
\usepackage{array}

% 周记抬头：\weekhead{周次}{时间区间}{本周主题}
\newcommand{\weekhead}[3]{%
\begin{center}
{\zihao{3}\bfseries 实习周记（#1）}\\[10pt]
{\zihao{-4}姓名：【姓名】\qquad 学号：【学号】}\\[4pt]
{\zihao{-4}时间：#2}\\[4pt]
{\zihao{-4}主题：#3}
\end{center}
\vspace{2pt}
\hrule height 0.8pt
\vspace{10pt}
\zihao{-4}%
}


\begin{document}

\weekhead{第 2 周}{2026 年 8 月 10 日至 8 月 14 日}{主网 trace 数据采集、访问画像与在线推理可行性分析}

本周的主线是拿到真实数据，验证想法有没有成立的前提。周一先把本地全节点切到离线模式跑起来，随后从高度 20,000,000 开始连续回放 1,001 个区块。采集方式是逐笔调用本地 Erigon 的 debug\_traceTransaction 接口，同时为了拿到训练需要的字段，在客户端的 trace 路径里埋了记录点，把 calldata 前几个参数、合约代码哈希和真实访问的槽位列表一起写盘，最终得到一份约 238 MB 的 JSONL，共 9.4 万笔有效交易。中间发现有两个区块过滤后只剩不含存储访问的交易，实际参与后续实验的是 999 个块。采集的流程本身不复杂，但全节点的数据目录有 1.7 TB，回放很慢，字段格式来回调又重跑了几轮，这类等待相当磨人，后来我学会先把解析和校验逻辑写好，再整批重跑。

数据到手后先做访问画像。统计结果是：每笔交易平均访问 16.2 个存储槽，其中 12.3 个是块内首次访问；全部访问事件中冷访问占 76%；87.9\% 的槽位在整块内只被一笔交易访问。写第一版统计小结时，我把 76\% 直接写成"76\% 的访问是冷的"，导师连着问了几个问题：分母是所有访问事件还是槽位个数，有没有按交易量加权，这个比例能不能推出大多数槽位是冷的。我回去核对之后才发现，自己最初的口径其实没有说清楚。后来我把口径固定为"冷访问事件占全部存储访问事件的比例"，并在论文的对应位置单独留了注释。这件事给我的教训是，一个百分比如果不写清分子和分母，作者容易含糊过去，读者更容易误解。

画像做完之后是两项分析。第一项按 (to, selector) 聚合，同一合约同一函数的交易，历史上访问过的槽位并集大多很小，83.5\% 的键并集在 50 个槽以内，少数并集很大的键承载了大量交易，呈明显的长尾分布。第二项是评估在线 ML 推理的可行性。我把 LightGBM 多标签模型在实验平台上测了单笔推理耗时，约 28 ms，而一次预取最多省 3 微秒，差了约三个数量级。为了把"推理也许能和执行重叠"的辩护也堵住，我按导师的建议引入覆盖参数 α，即使假设 α=1.0，即推理完全与执行并行，推理总耗时 14.78 秒，仍然远大于存储代价的节省 0.005 秒。结论是负面的，却直接决定了方案走向：ML 推理必须离开在线的关键路径。

于是方案从"交易前用模型快速推理"调整为 online/offline 分工：在线只做 O(1) 的哈希查表，模型在离线批量运行，把新发现的访问模式合并回表里。这个方向正好用上了并集分布的两条规律，并集小的键直接查表，并集大的键才值得交给模型精挑。周五和导师讨论后把方案定了下来，下周开始实现仿真器和离线增量管道。这两周让我体会到，数据先行不是一句口号，口径、分母这些细节不较真，后面每个结论都可能站不住。

\end{document}
